\subsubsection*{Solution}
The source-target convolution and its marginal saddle-point behavior are discussed by Hallatschek and Fisher. In particular, the convolution is dominated near $\tau=t/2$, and the shrinking saddle width produces a $\log t\log\log t$ correction to the leading $(\log t)^2$ behavior. \href{\detokenize{https://doi.org/10.1073/pnas.1404663111}}{PNAS article}, \href{\detokenize{https://www.rpgroup.caltech.edu/virus_bootcamp/assets/pdfs/pnas-hall-dsf-full.pdf}}{Supporting Information, Sec. SI2.C}.

Here, the exponent $\mu=2$ refers directly to the power of $\ell(t)$ in the stated equation.

\paragraph*{Step-by-Step Derivation}

We solve

\[
\int_0^t\ell(\tau)\ell(t-\tau)\,d\tau =K\ell^2(t). \tag{1}
\]

The logarithms imply that $t$ and $K$ are expressed in the same fixed time units.

\subparagraph*{1. Logarithmic variables}

It is convenient first to use natural logarithms:

\[
x=\ln t,\qquad L=\ln 2,\qquad F(x)=\ln \ell(e^x).
\]

Thus

\[
\varphi(z)=\frac{F(Lz)}{L}, \qquad z=\log_2t=\frac{x}{L}. \tag{2}
\]

With $\tau=tu$, equation (1) becomes

\[
t\int_0^1 \exp\!\left[F(x+\ln u)+F(x+\ln(1-u))\right]du =K e^{2F(x)}. \tag{3}
\]

Define

\[
H_x(u)=F(x+\ln u)+F(x+\ln(1-u)).
\]

By symmetry, its saddle lies at $u=1/2$. Set

\[
q=x-L,\qquad p=F'(q),\qquad r=F''(q).
\]

At the saddle,

\[
H_x\!\left(\frac12\right)=2F(q),
\]

and direct differentiation gives

\[
H_x''\!\left(\frac12\right)=8(r-p).
\]

Since asymptotically $p-r>0$, this is a maximum.
{\color{blue}
It is in fact the global maximum. This requires two statements: that $\tau=t/2$ is the unique maximum in the interior of the integration range, and that the endpoint regions are negligible.

\emph{Interior.} In terms of $\tau=tu$ the integrand is $\ell(\tau)\ell(t-\tau)$. Let $\gamma(\tau)=\frac{d}{d\tau}\ln\ell(\tau)=F'(\ln\tau)/\tau$ be the growth rate of the cluster. Then
\[
\frac{d}{d\tau}\ln\bigl[\ell(\tau)\ell(t-\tau)\bigr]=\gamma(\tau)-\gamma(t-\tau),
\qquad
\gamma'(\tau)=\frac{F''(\ln\tau)-F'(\ln\tau)}{\tau^{2}}<0 ,
\]
the inequality because $F'\gg F''$ (by (10) below, $F'(w)\simeq w/(2L)$ while $F''(w)\simeq1/(2L)$). So $\gamma$ is decreasing, and the derivative above is positive for $\tau<t/2$ (where $\tau<t-\tau$) and negative for $\tau>t/2$: wherever the asymptotic form of $F$ applies, the integrand rises until $\tau=t/2$ and falls afterwards.

\emph{Endpoint regions.} Near the endpoints ($\tau=O(1)$ or $t-\tau=O(1)$) the asymptotic form of $F$ does not apply and $\ell$ is given by its early-time behaviour instead. These regions contribute $O(\ell(t))$ to the left side of (1), smaller than the saddle-point contribution $\sim t\,\ell(t/2)^2$ by a factor $\exp[-ax^2+O(x\ln x)]$, hence negligible to all orders in $1/x$. This is also why all coefficients below are independent of the early-time behaviour of $\ell$; only the additive constant is not.
}
Introduce

\[
A=-H_x''\!\left(\frac12\right)=8(p-r).
\]

\subparagraph*{2. Saddle-point expansion}

The fourth derivative at the saddle is

\[
H_x^{(4)}\!\left(\frac12\right) = 32\left(F^{(4)}(q)-6F'''(q)+11F''(q)-6F'(q)\right).
\]

Because $F'(q)=O(x)$ while $F''(q)=O(1)$, the first non-Gaussian correction is

\[
\frac{H_x^{(4)}(1/2)}{8A^2} =-\frac{3}{8p} +O\!\left(\frac{\ln x}{x^2}\right).
\]

Consequently,

\[
\int_0^1e^{H_x(u)}du = e^{2F(q)} \sqrt{\frac{\pi}{4(p-r)}} \exp\!\left[ -\frac{3}{8p} +O\!\left(\frac{\ln^2x}{x^2}\right) \right]. \tag{4}
\]

Substitution into (3) and taking logarithms gives

\[
2\bigl[F(x)-F(x-L)\bigr] = x-\ln K+\frac12\ln\frac{\pi}{4} -\frac12\ln\!\bigl(F'(q)-F''(q)\bigr) -\frac{3}{8F'(q)} +O\!\left(\frac{\ln^2x}{x^2}\right). \tag{5}
\]

\subparagraph*{3. Asymptotic ansatz}

{\color{blue}
\emph{Origin of the ansatz.} The form of $F$ need not be guessed: it is generated iteratively by (5). Expanding the left side of (5) in $L$,
\[
2L F'(x)-L^2F''(x)+O(F''') = x-\ln K+\tfrac12\ln\tfrac{\pi}{4}-\tfrac12\ln\bigl(F'(q)-F''(q)\bigr)-\frac{3}{8F'(q)}+\dots
\]
Under the saddle-point premise $F''\ll F'$ on which (5) rests (verified a posteriori by (10)), the only possible leading balance is $2LF'=x$, giving $F=x^2/(4L)$. Inserting this on the right, $-\tfrac12\ln F'=-\tfrac12\ln x+O(1)$ generates, through $2LF'$, a term $\propto x\ln x$; the $O(1)$ constants generate a term $\propto x$; expanding $-\tfrac12\ln(F'-F'')$ and $-3/(8F')$ to relative order $1/x$ generates terms $(\ln x)/x$ and $1/x$ on the right, which integrate to $\ln^2x$ and $\ln x$. All further terms are $O(\ln^2x/x)\to0$. Since every step only produces powers of $x$ and $\ln x$, no terms of any other type can arise; carrying the iteration to $O(1)$ yields the ansatz below, which therefore contains every term that does not vanish as $x\to\infty$.
}
{\color{blue}Accordingly, take}

\[
F(x) = ax^2+bx\ln x+cx +d(\ln x)^2+e\ln x+O(1). \tag{6}
\]

Its relevant derivatives are

\[
F'(x) = 2ax+b(\ln x+1)+c +\frac{2d\ln x+e}{x} +O\!\left(\frac{\ln x}{x^2}\right),
\]

\[
F''(x) = 2a+\frac{b}{x} +O\!\left(\frac{\ln x}{x^2}\right).
\]

Expanding the left side of (5),

\begin{equation}
\begin{aligned}
2[F(x)-F(x-L)]
={}&4aLx+2bL\ln x
+2L(b+c)-2aL^2\\
&+\frac{4dL\ln x+2eL-bL^2}{x}
+O\!\left(\frac{\ln x}{x^2}\right).
\end{aligned}
\tag{7}
\end{equation}

On the right side,

\begin{equation}
\begin{aligned}
\ln\!\bigl(F'(q)-F''(q)\bigr)
={}&\ln x+\ln(2a)\\
&+\frac{1}{x}
\left[
\frac{b\ln x+b+c-2a}{2a}-L
\right]
+O\!\left(\frac{\ln^2x}{x^2}\right),
\end{aligned}
\tag{8}
\end{equation}

and

\[
\frac{1}{F'(q)} = \frac{1}{2ax} +O\!\left(\frac{\ln x}{x^2}\right). \tag{9}
\]

Matching successively the coefficients of $x$, $\ln x$, $1$, $(\ln x)/x$, and $1/x$ gives

\[
a=\frac{1}{4L}, \qquad b=-\frac{1}{4L}, \qquad d=\frac{1}{16L}, \tag{10}
\]

\[
c= \frac{ -\ln K+\frac12\ln\!\left(\frac{\pi L}{2}\right) +\frac12+\frac{L}{2} }{2L}, \tag{11}
\]

and

\[
e=-\frac{c}{2}+\frac{3}{8L}-\frac14. \tag{12}
\]

\subparagraph*{4. Conversion to binary logarithms}

Using $x=Lz$ and $\ln x=L\log_2z+\ln L$, equation (6) becomes

\[
\varphi(z) = \frac{z^2}{4} -\frac{z}{4}\log_2z +A_Kz +\frac{1}{16}\bigl(\log_2z\bigr)^2 +B_K\log_2z +O(1), \tag{13}
\]

where

\[
A_K = \frac14 \left[ 1+\log_2\!\left(\frac{\mathrm e\pi}{2K^2}\right) \right], \tag{14}
\]

and

\[
B_K = \frac18 \left[ 2\log_2\mathrm e-3 +\log_2\!\left(\frac{2K^2}{\pi}\right) \right]. \tag{15}
\]

The additive $O(1)$ term is not fixed because equation (1) is invariant under $\ell(t)\mapsto C\ell(t)$.

{\color{blue}
\paragraph*{Reviewer's verification}
The derivation above was checked in three independent ways.

\subparagraph*{(i) Symbolic re-derivation}
The order-by-order matching of Step 3 and the conversion of Step 4 were redone independently with SymPy and reproduce (10)--(12) and (14)--(15) exactly. Carrying the matching one order further gives the $O(1/x)$ corrections to (6), which are used in (iii).

\subparagraph*{(ii) Residual of the integral equation}
With (6) and (10)--(12) inserted into (3), and the $u$-integral evaluated numerically over $0.01<u<0.99$ (the rest is negligible, as discussed in Step 1), the residual $R(x)=x+\ln I-\ln K-2F(x)$ satisfies $xR\to0$ as $x=\ln t\to\infty$ ($xR=1.4\times10^{-6}$ at $x=10^7$ for $K=1$). An error $\delta$ in the coefficient $e$ of (6) would instead give $xR\to-2L\delta$, and an error in any of the other coefficients of (6) would make $xR$ diverge, so this confirms all five coefficients in (10)--(12).

\subparagraph*{(iii) Direct numerical solution}
Equation (1) was also solved numerically from explicit initial data up to $\ln t=300$: from $\ell(t)=1$ for $t\le1$ with $K=1$ and $K=3$, and from a different $O(1)$ early history with $K=1$. The difference between $\ln\ell$ and the expansion (6), including the $O(1/x)$ corrections from (i), settles to a constant: it varies by less than $5\times10^{-5}$ over $20\le\ln t\le300$ for $K=1$ (and by less than $2\times10^{-4}$ over $100\le\ln t\le300$ for $K=3$), whereas the $e\ln x$ term alone changes by $0.08$ (respectively $0.40$) over the same range. The constant depends on the initial data; nothing else does.

}
